% Source: https://artofproblemsolving.com/community/c3h2693136_vietas_on_double_roots
% Stable identity: aops-topic-2693136-q01
Vieta's formulas relate the coefficients to the roots \emph{counted with multiplicity}: a double root $r$ is listed twice. So let the roots be $r,r,s$. Then
\[2r+s=5,\qquad r^2+2rs=8,\qquad r^2s=4.\]
From the first, $s=5-2r$; substituting into the second gives $r^2+2r(5-2r)=8$, i.e. $3r^2-10r+8=0$, so $r=2$ or $r=\dfrac43$.
\begin{itemize}
\item $r=2$: $s=1$ and $r^2s=4$. This works.
\item $r=\dfrac43$: $s=\dfrac73$ and $r^2s=\dfrac{112}{27}\ne4$. This fails.
\end{itemize}
So the roots are $\boxed{2,2,1}$, i.e. $x^3-5x^2+8x-4=(x-2)^2(x-1)$.
